% 17 · Conclusions: the eight lessons, and what comes next
\section{CONCLUSIONS}
\label{sec:conclusions}

\deckslides{slides 48--49: the eight lessons and the take-away.}

This chapter collects the eight lessons the benchmark produced, each with the
number that supports it and the section that demonstrates the mechanism. They
are the same eight the lecture closes on; here they are expanded to the length
the evidence allows.

\subsection{Eight lessons from the benchmark}

\textbf{1 · All-sky collapses; cut a region.} One embedding of every clean star
in the survey buries each cluster in the field: the field is 357\,056 stars of
chemically smooth background, and 25 clusters of a few dozen members each are
not density peaks in that space. The published answer, and the one that works,
is to restrict the sky. \citet{Kos:17} used $40^\circ$ around the Pleiades, and
our region runs use a cone scaled to the cluster, $\max(3^\circ, 10\times$
diameter$)$. The cost of ignoring this is measurable: in the uniform all-sky
run the abundance-only t-SNE arm reaches recall $0.201$ at precision $0.240$
(\textbf{Table~\ref{tab:field}}), while the same pipeline in region mode on the
DR17 data reached recall $0.73$ at precision $0.47$
(\textbf{Table~\ref{tab:curse}}). Same methods, same clusters; the difference
is the contrast between the cluster and its immediate neighbourhood.

\textbf{2 · Precision is the hard part.} Read across the region-by-region sweep
and the macro precision is $0.28$ for t-SNE, $0.08$ for UMAP and $0.01$ for
EVoC (25 clusters, DR19). The field is chemically similar to the clusters (that
is what it means for a cluster to be dissolved in it) so a method that recovers
members also recovers their look-alikes. We can now put a number on the
mechanism that \citet{Casamiquela:21} report: they find that more than 70\,\%
of the groups in their best-case solution are statistical mixtures, and our own
cluster-only runs give $64$--$72\,\%$ of groups that recover no cluster at all
(\texttt{stat40} in \S\ref{sec:literature}). Their negative result is about
cluster-to-cluster separation and ours about field contamination, but both say
the same thing: at survey precision, the chemical differences between clusters
of the same age are comparable to the differences between a cluster and the
field.

\textbf{3 · Recall $\approx 1$ is usually a blob.} HDBSCAN$^*$ on the abundance
space swallows the dense field into a single cluster, after which any
best-overlap scorer reports that the true members are inside it: recall $1.000$
at precision $0.0013$ (UMAP on abundances, \textbf{Table~\ref{tab:field}}, and
note that this run has row-normalisation \emph{on} (it is the configuration
default, and it does not prevent the collapse). Two defences are built into
this workbook: report kNN purity, which cannot be gamed by a group's size; and
never quote recall without its precision. The third defence we tried)
row-normalisation. Is the subject of the quiet lesson below.

\textbf{4 · Globulars tag cleanly; open clusters do not.} The kNN purity of the
metal-poor globulars reaches $0.4$--$0.5$ while solar-metallicity open clusters
stay below $0.2$. Three effects push the same way: globulars are older, more
metal-poor, and formed from gas that no longer resembles the present-day
interstellar medium, so their chemistry is distinct from the field they sit in.
The Pleiades are the limiting case ($100$ Myr old, formed from gas the field
still resembles) and the prediction that chemistry fails there is part of the
assignment (\S\ref{sec:assignment}), to be reported as a null result when it
does.

\textbf{5 · t-SNE isolates, UMAP over-merges, EVoC is fast but coarse.} Three
sets of assumptions, three answers, and no winner. t-SNE produces the most
separated picture and is the least stable under row order
($0.218$--$0.560$ on the abundance matrix, \S\ref{sec:validation}). UMAP is
the fastest and the only arm we run on the full field, and it is the arm that
produces blobs. EVoC needs no clustering choice from the user, works on the
cosine geometry the data actually have, and has a seed spread
($\pm0.043$ on abundances) that overlaps the gaps being claimed between
methods.

\textbf{6 · Levers are fragile, and read them in pairs.} On the DR17 data,
weighting each element by $1/\sigma$ \emph{after} standardising lifted M~67's
t-SNE recall from $0.09$ to $0.44$: the largest single effect in the sweep, and
the fix for a bug that had made the same switch a no-op. On the DR19 data the
same switch moves recall from $0.10$ to $0.96$ while precision falls from
$0.14$ to $0.05$: t-SNE's cluster grows until it swallows the field. The region
sweep agrees, with precision falling $0.83 \rightarrow 0.24$. A lever selected
on one metric, one cluster and one seed is a hypothesis; it becomes a result
only when it is re-run on the current data with both halves of the metric pair
reported.

\textbf{7 · Spectra beat abundances, and not PCA.} A self-supervised masked
autoencoder on the raw spectrum separates the 25 clusters better than the 16
ASPCAP abundances ($0.760$ vs $0.578$ homogeneity, UMAP, the same 982 stars),
and a plain PCA of the same spectra does about as well ($0.754$). Two very
different representations of the same pixels land on the same score, which
tells you where the limit is: not in the model, but in what the spectra and the
metric jointly permit. On field retrieval the two tie to the third decimal
($0.418$/$0.619$ against $0.417$/$0.631$), and the only field-retrieval claim
we made for the spectral latent has been withdrawn for the reason given in
\S\ref{sec:provenance}.

\textbf{8 · Two surveys, two parameters.} APOGEE sees giants: the red clump
gives distances, and NGC~6819 comes out within $0.01$ mag of the literature
value (\textbf{Table~\ref{tab:dm}}). GALAH sees the main sequence and turnoff:
that is where age sensitivity lives. The joint fit is not there yet (the age
posterior ($\sigma_{\log a} \approx 0.97$) is close to the width of the flat
prior, so the age half of the promise needs a better likelihood before it shows
anything) but the distance half works, and the NGC~2243 experiment
(\S\ref{sec:physics}) shows what the combination can do when the distance is
pinned and the main sequence is added: an age residual improving from $0.52$ to
$0.09$.

\subsection{The quiet lesson}

We used to say that row-normalisation breaks the blob, with a precision gain
of three- to fivefold. Measured across the 25 clusters of the DR19 region sweep,
it does not: turning normalisation \emph{off} raises UMAP's mean precision by
$1.7\times$ ($0.08 \rightarrow 0.14$) and cuts the number of blobs from twelve to
five, while t-SNE is unchanged ($0.28 \rightarrow 0.26$) and EVoC moves by a
few thousandths. On M~67 the effect runs the other way: with normalisation on,
t-SNE precision is $0.83$; off, it is $0.28$. The lever is real, but its
direction depends on the method and the cluster, which is exactly what a
lever selected on one cluster at one seed looks like when it is re-measured
properly (\S\ref{sec:validation}). What survives as a rule is smaller and
duller: the precision--recall trade is there in every setting, and it is a
choice about which error you can live with, made explicit.

\begin{summary}[SUMMARY POINTS]
\begin{enumerate}
\item Clustering is a decision problem: which space, which metric, which notion
of ``cluster'', which scale, and every choice belongs in the write-up.
\item Geometry beats algorithm. The largest improvements in this project came
  from changing the representation (spectra instead of abundances,
  row-normalisation, 64 components instead of 256), not from changing the
  clusterer.
\item Density-based clustering needs a hierarchy or a persistence criterion.
  One global threshold cannot serve a field of 357\,056 stars and a cluster of
  6 members at the same time.
\item Embeddings are for neighbourhoods, not distances. A t-SNE or UMAP map is
  evidence about who is next to whom, and about nothing else.
\item Validate with something the algorithm never saw, and print the
  degenerate cases. Largest-group fraction, number of groups, both halves of
  the metric pair.
\item Measure the ceiling and report the gap. Kinematics are not a competitor;
  they are the ceiling that chemistry is being compared with.
\item A lever is a hypothesis. Re-run it on the current data, with both
  halves of the metric pair, before quoting it.
\item Membership quality propagates into physics: M~67's isochrone age moves
  from $1.89$ to $2.65$ Gyr, against a literature value of $2.82$, as the
  membership list improves.
\end{enumerate}
\end{summary}

\begin{issues}[FUTURE ISSUES]
\begin{enumerate}
\item \textbf{Re-embed on a single data product.} The mixed-product latent of
  \S\ref{sec:provenance} invalidates every field-retrieval number computed from
  it. Uniform \texttt{mwmStar} spectra for the backfilled members are already
  downloaded ($736$ spectra), and the re-embedding run awaits the model
  checkpoint. Until it is done, the field-retrieval rows for the spectral
  latent should be read as provisional.
\item \textbf{Collapse the duplicate rows.} 298 member rows share an
  \texttt{APOGEE\_ID} with another row. Until the sample is deduplicated, every
  macro average over clusters is weighted by something that is not a star.
\item \textbf{Benchmark against \citet{Spina:25}.} \S\ref{sec:literature}
covers the published negative result (their clustering step on our stars, their
metric triple on our scores) but a common-population comparison with the
graph-attention autoencoder, on the same clusters and the same two tasks, is
still the experiment this work is missing.
\item \textbf{Test the strong-tagging claim.} Everything here is measured with
  a kinematic prior available. Recovering a dispersed cluster from the field
  alone remains the open problem, and the honest current status is
  precision $\approx 0.62$ at recall $0.42$.
\item \textbf{A metric for high-dimensional chemical space.} The dilution
  result of \S\ref{sec:benchmark} is a statement about Euclidean distance, not
  about chemistry. A metric that respects the compositional structure of
  abundance vectors might make concatenation behave, and nobody has
  demonstrated one.
\end{enumerate}
\end{issues}

\begin{issues}[DISCLOSURE AND PROVENANCE]
Every number in this workbook is produced by code in the accompanying
repository at a recorded commit, and the scripts that generate the tables are
named in the sections that use them. Three things are worth distinguishing.
(1) The duplicate-row and row-order audit of \S\ref{sec:validation}, the
product-mismatch diagnosis of \S\ref{sec:provenance}, and the withdrawal of
the field-retrieval claim for the mixed latent are the repository's own
findings, reported here as they stand. (2) The K-means row of
\textbf{Table~\ref{tab:honest}} was added for this workbook
(\texttt{article/scripts/kmeans\_baseline.py}) and had not been measured
before. (3) The figures are the lecture's, reused from the deck so that the
workbook and the slides describe the same drawings; the animations are
flattened to filmstrips, and the deck itself is the primary reference for
anything that moves.
\end{issues}

\subsection{Acknowledgments}

The workbook accompanies the 2026 IAA--SO school on advanced neural networks,
and is built on the accompanying repository's benchmark. The data are the
public work of the SDSS-V/APOGEE, GALAH and Gaia collaborations, and the
methods are the work of the authors cited throughout. The teaching figures and
the dataset design come from the lecture; the isochrone fits use the PARSEC
grids through AStECA and \texttt{emcee}; the pipeline is written on
\texttt{numpy}, \texttt{pandas}, \texttt{scikit-learn}, \texttt{astropy},
\texttt{matplotlib} and \texttt{umap-learn}, with EVoC for the production
runs. Errors of interpretation that remain are ours, and the exercises exist
so that the next pass is yours.

\begin{figure}[htbp]\centering
  \begin{tabular}{@{}c@{\hskip14pt}c@{}}
    \fig[0.22]{qr_school.png} &
    \fig[0.22]{qr_library.png}
  \end{tabular}
  \caption{The school programme and the reference library. The deck itself,
with its animations, is at the URL on the title page and in the QR code of
\textbf{Figure~\ref{fig:qrcodes}}.}
  \label{fig:qrclosing}
\end{figure}
